\documentclass[11pt,a4paper]{article}
\usepackage[utf8]{inputenc}
\usepackage[T1]{fontenc}
\usepackage{estilo_capstone}

\titulo{Manual técnico y de despliegue — Evento Cultural}
\subtitulo{Manual técnico/despliegue (artefacto 8 del marco ágil, anexo metodológico)}
\autor{Roberto Riffo (frontend) · José Miguel Moncada (BD/QA) · David Herrera (backend e integraciones)}
\fecha{25 de septiembre de 2026}

\newcolumntype{L}[1]{>{\raggedright\arraybackslash}p{#1}}

\begin{document}
\maketitlecapstone

\section{Composición del sistema}

\begin{table}[h!]
\centering\small
\begin{tabular}{@{}L{3.2cm}L{5.4cm}L{6.2cm}@{}}
\toprule
\textbf{Parte} & \textbf{Tecnología} & \textbf{Rol} \\
\midrule
Scraper & Python (conectores, NVIDIA NIM, checkpoints SQLite) & Recolección, normalización y enriquecimiento de eventos desde 21 fuentes públicas \\
Base de datos & Supabase (PostgreSQL + RLS) & Catálogo canónico, bitácora de ejecuciones y evidencia de extracción \\
Frontend & Angular 22 + Ionic 9 (web y móvil) & Descubrimiento: búsqueda, filtros, mapa, calendario y recomendaciones \\
Automatización & GitHub Actions & Ejecución periódica/manual del scraper con pruebas \\
\bottomrule
\end{tabular}
\caption{Componentes del sistema.}
\end{table}

\section{Ejecución local}

\subsection{Scraper (Windows)}

\begin{verbatim}
instalar.bat        REM crea venv, instala dependencias y Chromium de Playwright
ejecutar.bat        REM pipeline completo (extraccion + exportacion + sincronizacion opcional)
python -m unittest discover -s tests -v
\end{verbatim}

Las variables de entorno se documentan en \code{.env.example} del proyecto;
incluyen credenciales de Supabase y el endpoint de NVIDIA NIM. \textbf{No}
versionar secretos: el workflow los consume como secretos de GitHub.

\subsection{Frontend}

\begin{verbatim}
npm install
npm start           # ng serve en http://localhost:8100
ng test             # pruebas (Vitest)
npx ng build --configuration production
\end{verbatim}

En Windows existen además \code{install.bat} y \code{run.bat} para el entorno
local. El frontend funciona con catálogo mock en memoria; la conexión a
Supabase es un pendiente declarado.

\section{Estrategia de despliegue (sin Docker)}

El equipo decidió \textbf{no usar Docker}; el anexo metodológico se cumple
documentando esta decisión y el estado real. El despliegue se compone así:

El sistema no se aloja en contenedores: las plataformas de alojamiento y
automatización son GitHub Actions, Vercel y Capacitor. El estado de cada
componente:

\begin{table}[h!]
\centering\small
\begin{tabular}{@{}L{3.0cm}L{5.6cm}L{6.2cm}@{}}
\toprule
\textbf{Componente} & \textbf{Plataforma} & \textbf{Estado} \\
\midrule
Scraper & GitHub Actions (workflow \code{scrape-events.yml}) & \textbf{Desplegado, con ejecución manual y cron diario activo} desde el 24-09 (00:37, America/Santiago; corridas \#9--\#11 verificadas del 25 al 27-09, evidencia en Evidencias Proyecto); la corrida del 27-09 falló en pruebas (Plan de pruebas §5.1) \\
Web & Vercel & La página se alojará en Vercel: despliegue en preparación \\
Móvil (Android) & Capacitor 8 + Android Studio & La misma web se convertirá en aplicación mediante Capacitor 8: conversión en preparación (dependencias presentes en \code{package.json}) \\
\bottomrule
\end{tabular}
\caption{Estrategia de despliegue por componente y su estado real.}
\end{table}

\subsection{Scraper en GitHub Actions (desplegado, en testeo)}

\begin{itemize}
  \item El workflow está desplegado y el equipo se encuentra probando el
        despliegue: correr el pipeline con pruebas y sincronización a
        Supabase, usando secretos de repositorio (nunca credenciales en el
        código).
  \item Evidencia de ejecución: captura \code{Evidencia\_ejecución\_scraper\_en
        \_GitHub\_Actions.PNG} en Evidencias Proyecto.
  \item La periodicidad ya está definida y activa: \code{schedule} con cron diario
        a las 00:37 (America/Santiago) desde el 24-09. Para dar el despliegue por
        completado, el equipo debe corregir la prueba dependiente del tiempo del 27-09
        (Plan de pruebas §5.1) y obtener una corrida programada completa en verde.
\end{itemize}

\subsection{Web en Vercel (alojamiento en preparación)}

\begin{enumerate}
  \item Conectar el repositorio del frontend a Vercel (framework Angular/Ionic
        detectado automáticamente o con \code{ng build --configuration
        production} como comando de build).
  \item Configurar variables de entorno de Supabase en el proyecto de Vercel
        (mismas claves que \code{.env.example}).
  \item Publicar y registrar la URL como evidencia de despliegue.
\end{enumerate}

\subsection{De la web a la app: Capacitor (conversión en preparación)}

La aplicación Android es la misma web envuelta con Capacitor; pasos de la
conversión:

\begin{enumerate}
  \item \code{npx cap add android} (si la plataforma no existe) y
        \code{npx cap sync}.
  \item Abrir en Android Studio, compilar el APK/AAB de prueba y registrar la
        evidencia de la compilación.
\end{enumerate}

\section{Variables de entorno}

\begin{table}[h!]
\centering\small
\begin{tabular}{@{}L{5.2cm}L{9.6cm}@{}}
\toprule
\textbf{Grupo} & \textbf{Uso} \\
\midrule
Supabase (URL y clave de servicio) & Sincronización del scraper y, en el futuro, del frontend \\
NVIDIA NIM (endpoint y clave) & Enriquecimiento de eventos en el pipeline \\
\bottomrule
\end{tabular}
\caption{Grupos de variables de entorno (nombres exactos en \code{.env.example}).}
\end{table}

\section{Solución de problemas comunes}

\begin{itemize}
  \item \textbf{Pipeline detenido a mitad:} los checkpoints en SQLite permiten
        reanudar; volver a ejecutar \code{ejecutar.bat}.
  \item \textbf{Consola del navegador con rechazos de View Transitions:} los
        rechazos benignos ya se suprimen en \code{main.ts}; cualquier otro
        error de consola se trata como falla de la verificación.
  \item \textbf{Puerto 8100 ocupado:} liberar el proceso de \code{ng serve}
        previo.
\end{itemize}

\end{document}
